\section{Semiample models and rounded section spaces}
\label{sec:reductions}

We pass from the original pair to a semiample model while preserving
every rounded section space. The following comparison keeps every
integer degree, even when the corresponding adjoint multiple is not
Cartier.

We use the following precise theorem of the companion
\emph{Lifting sections from the reduced support of an adjoint}
\cite[Theorem~\BANumber]{OpenAILifting}.

\begin{theorem}[Abundance after nonvanishing]\label{thm:fourfold}
Let $(X,\Delta)$ be a projective lc rational pair over $\C$ of
 dimension at most four, with effective boundary and nef
 $\Q$-Cartier adjoint $D=K_X+\Delta$. If $\kappa(X,D)\ge0$,
 then $D$ is semiample. If $\kappa(X,D)=0$, then $D\sim_{\Q}0$.
\end{theorem}

This is the companion's independent after-nonvanishing theorem. Its proof
uses the supported-boundary theorem and lower-dimensional abundance;
it does not use the separate fourfold nonvanishing companion. The boundary
coefficients may be arbitrary rational numbers in $[0,1]$. The conclusion
is qualitative: it supplies neither a uniform index nor a uniform
section degree. If a generated Cartier multiple has Iitaka dimension
zero, its morphism has image a point and a generating section is nowhere
zero, which explains the actual rational linear triviality.

For the original pair we first run an adjoint-negative program and
compare all section spaces. For the moving-divisor application in
Lemma~\ref{lem:index-moving-divisor}, the same procedure applies to a
small rational perturbation $(V,\delta D)$. There it supplies only a
contraction. Uniform statements enter after the boundary is restored
to zero and canonical crepancy is proved separately.

\begin{proposition}\label{prop:semiample-model}
Let $(X,\Delta)$ be a projective lc rational pair of dimension four
over $\C$, with $D=K_X+\Delta$ \(\Q\)-Cartier and
$\kappa(X,D)\geq0$. There is a projective \(\Q\)-factorial dlt pair
$(X',\Delta')$ with semiample adjoint $D'=K_{X'}+\Delta'$ such that,
under the identification of rational function fields and compatible
canonical divisors,
\begin{equation}\label{eq:all-rounded-sections}
 H^0\bigl(X,\OO_X(\lfloor mD\rfloor)\bigr)
 =H^0\bigl(X',\OO_{X'}(\lfloor mD'\rfloor)\bigr)
 \qquad(m\geq1).
\end{equation}
The coefficients of $\Delta'$ belong to those of $\Delta$ together
with $1$. The map is obtained by a crepant dlt modification followed
by adjoint-negative birational steps. For a \(\Q\)-factorial klt
input the latter steps can be taken directly, with klt outputs.
The semiample contraction $f:X'\to Z$ has
$\dim Z=\kappa(X,D)$ and
\[
 D'\sim_{\Q}f^*A
\]
for an ample rational Cartier divisor $A$ on the normal projective
variety $Z$.
\end{proposition}

\begin{proof}
We first establish the section comparison used throughout the
construction. For a nonzero rational function $v$ and a rational
divisor $L$ on a normal variety,
\begin{equation}\label{eq:floor-pole-test}
 v\in H^0\bigl(\OO(\lfloor mL\rfloor)\bigr)
 \quad\Longleftrightarrow\quad
 \Div(v)+mL\geq0.
\end{equation}
The equivalence holds because the coefficients of $\Div(v)$ are
integers. Effective \(\Q\)-Cartier pullback and pushforward give
both inclusions for any crepant birational morphism. More generally,
suppose that two models have a common resolution with maps $p,q$
to their input and output, and that their adjoints satisfy
\[
 p^*D_{\mathrm{in}}=q^*D_{\mathrm{out}}+E,
 \qquad E\geq0,\qquad E\text{ is }q\text{-exceptional}.
\]
A section on the output pulls back, remains effective after adding
$mE$, and pushes down to a section on the input. A section on the
input pulls back and pushes by $q$ to a section on the output, since
$q_*E=0$. Applying \eqref{eq:floor-pole-test} proves equality for
every $m$, without assuming that $mD$ is Cartier. Thus the equality
survives any sequence of these comparisons.

We now construct the models to which this comparison applies.
Take a crepant projective \(\Q\)-factorial dlt modification. This is
the usual dlt modification theorem; see, for example,
\cite[Theorem~2.8]{FujinoGongyo2014}. Its boundary consists of the
strict transform of $\Delta$ and exceptional components of coefficient
one.

Run an adjoint-negative minimal model program. We recall the existence
and termination inputs needed here. On a fixed negative extremal ray,
slightly decreasing the coefficient-one part gives a klt boundary
for which that ray is still negative. Klt contraction and flip
existence apply \cite{BCHM2010}. To identify its flip with that of the
original adjoint, write the perturbed klt adjoint as $D_\epsilon$ and
the contraction as $g$. Relative Picard number one gives a positive
rational number $a$ such that $L=D-aD_\epsilon$ is numerically trivial
over $g$. For sufficiently divisible $k$, both $kL$ and $-kL$ are
Cartier, and their differences from $D_\epsilon$ are relatively ample.
The relative base point free theorem makes both relatively semiample.
Their associated maps are constant on each connected fibre of $g$,
so a further multiple of $L$ descends to a Cartier divisor on the
normal contraction base. Thus $L$ is a rational pullback. On the
small flipped model, the transform of $D$ is consequently a positive
rational multiple of the relatively ample transform of $D_\epsilon$
plus that pullback. It is relatively ample as required. The original
dlt condition is preserved by these negative steps. A fibre-type end
is excluded by \(\Q\)-linear effectivity of the adjoint, which follows
from $\kappa\geq0$. Indeed, a curve in a general fibre not contained in the
support of an effective representative cannot have negative
intersection with that representative. Divisorial contractions
decrease the Picard number, and the remaining flips terminate by
\cite[Theorem~1.1]{ChenTsakanikas2023}, applied with zero nef
data to the pseudo-effective rational lc adjoint. The output is
therefore nef and \(\Q\)-factorial dlt. Starting with a klt pair
keeps all these outputs klt.

The initial modification is crepant, and discrepancy comparison for
each adjoint-negative step has exactly the effective exceptional form
used above. Iteration proves \eqref{eq:all-rounded-sections}; in
particular the nef output still has nonnegative Kodaira dimension.
Theorem~\ref{thm:fourfold} makes it semiample. A generated Cartier
multiple, followed by Stein factorization, gives the stated
contraction and ample rational divisor.

We also identify its field with the whole embedded Iitaka field.
For sufficiently divisible $a$, the divisor $aD'$ differs from $f^*(aA)$
by an integral principal divisor. Complete systems of sufficiently high
multiples of the ample divisor $aA$ generate $\C(Z)$, so
$\C(Z)\subset K(D')$. Conversely, take two nonzero sections in any
$V_m(D')$. Raising both to a sufficiently divisible common power puts
them in such a pullback degree. Their ratio to that power lies in
$\C(Z)$ by $f_*\OO_{X'}=\OO_Z$. The original ratio is algebraic over
$\C(Z)$, which is relatively algebraically closed in $\C(X')$ because
$f$ is a contraction. Thus every rounded-degree ratio already lies in
$\C(Z)$. Together with \eqref{eq:all-rounded-sections}, this proves
$K(D)=K(D')=\C(Z)$ inside the common function field.
\end{proof}

\begin{lemma}\label{lem:multiples-full-field}
Let $D$ be a rational Cartier divisor on a normal projective variety
such that $D\sim_{\Q}f^*A$ for a contraction $f:X\to Z$ and an
ample rational Cartier divisor $A$. If one nonempty complete system
$|\lfloor mD\rfloor|$ generates $\C(Z)$ inside $\C(X)$, then every
positive integer multiple of $m$ does so as well.
\end{lemma}

\begin{proof}
The field identification just proved uses only $D\sim_\Q f^*A$
and the contraction property, so every rounded-degree ratio lies in
$\C(Z)$. Fix a nonzero section $s$ in degree $m$. For each $k\ge1$,
multiplication by $s^{k-1}$ embeds $V_m(D)$ in $V_{km}(D)$: the
pole inequalities add, without a Cartier assumption. Ratios of the
resulting sections are unchanged. Hence the degree-$km$ ratios contain
the whole field generated in degree $m$, and cannot generate anything
beyond $\C(Z)$.
\end{proof}
